\section{Implementación}

\subsection{Visión general}

La implementación del analizador sintáctico se compone de cuatro archivos fuente en C y dos archivos de especificación de herramientas (Flex y yacc/Bison), todos ellos completamente comentados conforme al estándar de documentación del curso.

\begin{table}[H]
\centering
\caption{Archivos fuente del analizador sintáctico y su función.}
\label{tab:archivos}
\begin{tabular}{@{}llp{6.5cm}@{}}
\toprule
\textbf{Archivo} & \textbf{Herramienta} & \textbf{Función} \\
\midrule
\texttt{parser/parser.y} & yacc/Bison & CFG, acciones semánticas, gestión de errores \\
\texttt{parser/symtab.c} & C & Implementación de la tabla de símbolos del parser \\
\texttt{parser/symtab.h} & C & Interfaz (tipos y prototipos) de \texttt{symtab.c} \\
\texttt{src/driver.c}    & C & Punto de entrada; modos \texttt{-t} y parser completo \\
\texttt{include/tokens.h}& C & Constantes numéricas de tokens compartidas \\
\texttt{lexer/lexico\_equipo\_11\_scanner.l} & Flex & Scanner; genera \texttt{INDENT}/\texttt{DEDENT} \\
\bottomrule
\end{tabular}
\end{table}

\subsection{Estructura del proyecto}
\begin{verbatim}
analizador-sintactico/
  include/tokens.h             -- IDs de token documentados (Cuadro 1)
  lexer/lexico_equipo_11_scanner.l  -- Scanner Flex con indentacion Python
  lexer/lexer.h                -- API publica del scanner
  parser/parser.y              -- CFG + acciones semanticas (yacc/Bison)
  parser/symtab.c / symtab.h   -- Tabla de simbolos del parser
  src/driver.c                 -- Driver: modos -t (lexer) y parser
  build/triton_parser          -- Ejecutable final
  tests/{valid,invalid}/       -- Casos de prueba del equipo
  results/{parser,jsonl,benchmark}/  -- Resultados y logs
\end{verbatim}

\subsection{Compilación y ejecución}

\begin{verbatim}
git clone https://github.com/G4ndhiV/Analizador-sint-ctico.git
cd Analizador-sint-ctico
make build               # compila triton_parser
make test                # 12 casos unitarios
make jsonl-test          # 200 casos JSONL (curated + adversarial)
make benchmark           # 70 kernels benchmark/kernels/
./build/triton_parser archivo.tri        # modo parser
./build/triton_parser -t archivo.tri     # modo solo lexico
\end{verbatim}

% ─────────────────────────────────────────────────────────────────────────────
\subsection{Código fuente completo y comentado}

En esta sección se presenta el código fuente completo del analizador sintáctico.
Cada archivo incluye:
\begin{enumerate}
  \item Cabecera de archivo con propósito, descripción y relación con otros módulos.
  \item Comentario de cada función/sección explicando el \textit{qué}, \textit{por qué} e invariantes no evidentes.
  \item Comentarios inline en la lógica no trivial (manejo de indentación, desambiguación de operadores, diferimiento de tokens).
\end{enumerate}

% ─── parser/symtab.h ────────────────────────────────────────────────────────
\subsubsection{Archivo: \texttt{parser/symtab.h}}

Este archivo define los cuatro tipos de entrada de la tabla de símbolos
(\texttt{FuncEntry}, \texttt{ParamEntry}, \texttt{VarEntry}, \texttt{ImportEntry})
y declara la API que las acciones semánticas de \texttt{parser.y} invocan
durante el análisis. La trazabilidad directa entre los entregables del análisis
(funciones, parámetros, variables, imports) y los \textit{structs} de este archivo
satisface el Entregable~2 del reto.

\lstinputlisting[language=C, caption={\texttt{parser/symtab.h}: tipos e interfaz de la tabla de símbolos.}]{../parser/symtab.h}

% ─── parser/symtab.c ────────────────────────────────────────────────────────
\subsubsection{Archivo: \texttt{parser/symtab.c}}

\texttt{symtab.c} implementa las cuatro tablas estáticas de la tabla de símbolos
del parser. Cada función de inserción es invocada desde una acción semántica
específica de \texttt{parser.y}:

\begin{itemize}
  \item \texttt{symtab\_add\_import()} $\leftarrow$ acción de \texttt{import\_stmt}.
  \item \texttt{symtab\_add\_function()} $\leftarrow$ acción de \texttt{def\_stmt} (tras el~\texttt{':'}).
  \item \texttt{symtab\_add\_param()} $\leftarrow$ acción de \texttt{param}.
  \item \texttt{symtab\_add\_variable()} $\leftarrow$ acción de \texttt{assign\_stmt} e \texttt{ident\_stmt}.
\end{itemize}

La función \texttt{copy\_str()} garantiza que ningún campo de los \textit{structs}
desborde su tamaño fijo. La detección de duplicados en \texttt{symtab\_add\_variable()}
asegura que la misma variable no aparezca dos veces en la misma tabla de símbolos.

\lstinputlisting[language=C, caption={\texttt{parser/symtab.c}: implementación de la tabla de símbolos.}]{../parser/symtab.c}

% ─── src/driver.c ───────────────────────────────────────────────────────────
\subsubsection{Archivo: \texttt{src/driver.c}}

El \textit{driver} es el punto de entrada del ejecutable \texttt{triton\_parser}.
Implementa dos flujos de ejecución:

\begin{description}
  \item[Modo parser (por defecto)] Activa \texttt{parser\_mode = 1}, llama a
    \texttt{lexer\_reset\_for\_parse()} y \texttt{symtab\_reset()}, invoca
    \texttt{yyparse()} y verifica la ausencia de tokens sobrantes.
    Si todo es correcto imprime \texttt{PARSE\_OK} seguido de las dos tablas
    de símbolos (parser y léxica).
  \item[Modo solo-léxico (\texttt{-t})] Activa \texttt{lexer\_debug\_mode = 1}
    y consume tokens con un bucle \texttt{while(yylex())}, reproduciendo la
    salida del reporte léxico (casos TC-01\,--\,TC-04).
\end{description}

\begin{description}
  \item[Código de salida 0] Análisis exitoso.
  \item[Código de salida 1] Error sintáctico (\texttt{syntax\_error\_seen}).
  \item[Código de salida 2] Error léxico (\texttt{lexical\_error\_seen}).
\end{description}

\lstinputlisting[language=C, caption={\texttt{src/driver.c}: punto de entrada del analizador sintáctico.}]{../src/driver.c}

% ─── parser/parser.y ────────────────────────────────────────────────────────
\subsubsection{Archivo: \texttt{parser/parser.y}}

Este es el núcleo del analizador sintáctico. El archivo se divide en cuatro
secciones yacc:

\begin{enumerate}
  \item \textbf{Preámbulo C} (\texttt{\%\{...\%\}}): variables globales de
    contexto (\texttt{current\_func}, \texttt{param\_pos\_counter},
    \texttt{triton\_decorator\_pending}), funciones auxiliares de recuperación de
    lexemas (\texttt{id\_text}, \texttt{int\_text}, \ldots) y prototipo de
    \texttt{validate\_stmt\_line\_end}.

  \item \textbf{Declaraciones} (\texttt{\%union}, \texttt{\%token}, \texttt{\%type},
    precedencias): el \texttt{\%union} define los dos tipos semánticos
    (\texttt{int i} para índices, \texttt{char *s} para lexemas); los tokens
    numéricos coinciden con el Cuadro~1 del reporte léxico; las directivas
    \texttt{\%left}/\texttt{\%nonassoc}/\texttt{\%right} resuelven conflictos
    de precedencia y el \textit{dangling-else}.

  \item \textbf{Reglas de la CFG} (\texttt{\%\%...}): 41 no-terminales y 105
    alternativas (P-001\,--\,P-105) documentadas en el Apéndice~\ref{ap:cfg-completa}.
    Las acciones semánticas insertan entradas en \texttt{symtab} y validan
    restricciones que la gramática por sí sola no puede expresar
    (\texttt{delim\_is}, \texttt{validate\_stmt\_line\_end}).

  \item \textbf{Funciones C} (\texttt{\%\%} final): \texttt{validate\_stmt\_line\_end}
    usa \texttt{yypeek()} para detectar múltiples sentencias en la misma línea;
    \texttt{yyerror} formatea los mensajes de error sintáctico en el formato
    estándar del curso.
\end{enumerate}

\paragraph{Trazabilidad CFG $\to$ código.}
Cada producción de la CFG resumida en la Sección~\ref{sec:cfg_resumen}
se implementa directamente en una regla de \texttt{parser.y}. La
correspondencia es uno-a-uno: no existe código de manejo sintáctico fuera
de este archivo (excepto la generación de INDENT/DEDENT en el lexer,
documentada en la Sección~\ref{sec:modificaciones_lexer}).

\lstinputlisting[language=C, caption={\texttt{parser/parser.y}: especificación yacc completa del analizador sintáctico.}, label={lst:parser-y-gramatica}]{../parser/parser.y}

% ─── Ejemplo de salida ──────────────────────────────────────────────────────
\subsection{Ejemplo de salida del parser (Entregable~5)}

Al analizar un kernel mínimo \texttt{triton\_kernel\_min.tri}:
\begin{verbatim}
PARSE_OK: archivo 'triton_kernel_min.tri' sintacticamente valido.

=== PARSER SYMBOL TABLE ===
--- IMPORTS ---
1  line=1  module=triton         alias=-
2  line=2  module=triton.language  alias=tl
--- FUNCTIONS ---
1  line=5  name=add_kernel  arity=5  triton_jit=1
--- PARAMETERS ---
1  line=5  func=add_kernel  param=x_ptr    pos=0  annot=-
2  line=5  func=add_kernel  param=y_ptr    pos=1  annot=-
3  line=5  func=add_kernel  param=output   pos=2  annot=-
4  line=5  func=add_kernel  param=n_elements  pos=3  annot=-
5  line=5  func=add_kernel  param=BLOCK_SIZE  pos=4  annot=-
--- VARIABLES ---
1  line=7  scope=add_kernel  name=pid
--- PARSER SYMBOL TABLE ---
=== SYMBOL TABLE: IDENTIFIERS ===
1    triton
2    tl
...
\end{verbatim}

\subsection{Mensajes de error (Entregable~4)}

Los mensajes de error están estandarizados en dos formatos, generados
respectivamente por el scanner y por \texttt{yyerror()}:

\begin{verbatim}
LEXICAL_ERROR:  token_id=<999> line=<n> lexeme="<lex>" (<descripcion>)
SYNTAX_ERROR:   line=<n> near='' message='<descripcion>'
\end{verbatim}

Ejemplos reales de la suite de pruebas:
\begin{verbatim}
LEXICAL_ERROR: token_id=<999> line=3 lexeme="x$" (identificador con caracteres invalidos)
SYNTAX_ERROR: line=5 near='' message='syntax error, unexpected TOKEN_NEWLINE'
SYNTAX_ERROR: line=3 near='' message='multiples sentencias en la misma linea'
SYNTAX_ERROR: line=4 near='indent' message='Indentacion inconsistente'
\end{verbatim}

\subsection{Explicación completa de las secciones de Flex (integración con yacc)}
\label{sec:flex_integracion}

En lenguajes con indentación significativa como Python/Triton, el analizador
sintáctico no puede operar de forma aislada: depende de que el analizador
léxico realice un trabajo \emph{sintético} avanzado. Al trasladar la gestión
de la pila de indentación (\texttt{INDENT}/\texttt{DEDENT}) y el control de
múltiples sentencias en una línea al archivo \texttt{.l}, el código de Flex
pasa a formar parte del diseño sintáctico. Esta subsección explica las
\textbf{tres secciones} del archivo
\texttt{lexer/lexico\_equipo\_11\_scanner.l} con foco en su integración con
\texttt{parser.y}; el scanner completo es la versión integrada del analizador
léxico del Equipo~11~[3].

\subsubsection{Sección 1 --- Declaraciones de Flex (\texttt{\%\{...\%\}} y definiciones)}

El preámbulo \texttt{\%\{...\%\}} establece el puente con Bison:

\begin{itemize}
  \item \texttt{\#include "y.tab.h"} hereda los \textbf{códigos de token} de
    Bison (\texttt{TOKEN\_IDENTIFIER}, \texttt{KW\_DEF}, \texttt{TOKEN\_INDENT},
    \texttt{TOKEN\_DEDENT}, \ldots) y el tipo \texttt{YYSTYPE} del
    \texttt{\%union}; \texttt{\#include "tokens.h"} aporta los IDs numéricos
    del Cuadro~1.
  \item \textbf{Estado compartido con el parser} declarado como variables de
    módulo: \texttt{parser\_mode}, \texttt{lexical\_error\_seen},
    \texttt{extern int syntax\_error\_seen}, la pila
    \texttt{indent\_stack[256]} con \texttt{indent\_depth}, la cola
    \texttt{pending\_tokens[]}, los contadores \texttt{paren\_depth}/%
    \texttt{bracket\_depth} (para no emitir \texttt{DEDENT} dentro de
    paréntesis) y el historial \texttt{lex\_prev}, \texttt{lex\_prev2},
    \texttt{lex\_prev\_line}, \texttt{saw\_assign\_this\_line} usado para
    detectar múltiples sentencias.
  \item \texttt{extern YYSTYPE yylval} es el canal por el que el scanner
    entrega valores semánticos (índices de tabla en \texttt{yylval.i},
    lexemas en \texttt{yylval.s}) al parser.
  \item \textbf{Mecanismo clave de integración}: la macro
    \texttt{YY\_DECL} renombra la función generada por Flex a
    \texttt{yyflexlex()}, y un \texttt{yylex()} escrito a mano la envuelve
    para \textbf{inyectar primero} los \texttt{INDENT}/\texttt{DEDENT}
    pendientes antes de devolver el token real a Bison.
\end{itemize}

\lstinputlisting[language=C, firstline=90, lastline=96, caption={Declaraciones de integración: \texttt{yylval} compartido y el truco \texttt{YY\_DECL} (\texttt{yyflexlex} envuelto por \texttt{yylex}).}]{../lexer/lexico_equipo_11_scanner.l}

Las definiciones de patrones relevantes para el parser son \texttt{ID}
(\texttt{[a-zA-Z\_][a-zA-Z0-9\_]*}), \texttt{BAD\_ID} (identificadores con
caracteres prohibidos), \texttt{WS\_COUNT} (\texttt{\textasciicircum[ \textbackslash t]+},
indentación al inicio de línea) y \texttt{NEWLINE} (\texttt{\textbackslash n+}).

\subsubsection{Sección 2 --- Reglas y patrones (\texttt{\%\% ... \%\%})}

Cada regla, en \texttt{parser\_mode}, fija \texttt{yylval} y ejecuta
\texttt{return TOKEN\_*}, entregando el token a \texttt{yyparse()}. Tres
reglas impactan directamente al análisis sintáctico:

\begin{itemize}
  \item \textbf{\texttt{WS\_COUNT}} (inicio de línea): calcula el ancho de
    indentación con \texttt{indent\_width()} y \texttt{indent\_line\_kind()} y
    llama a \texttt{handle\_indent()}, que compara contra la cima de la pila
    y emite \texttt{TOKEN\_INDENT} (sangría mayor) o uno o varios
    \texttt{TOKEN\_DEDENT} (sangría menor) que el parser consume en la regla
    \texttt{suite}.
  \item \textbf{\texttt{BAD\_ID}} (declarada \emph{antes} que \texttt{ID}
    para tener prioridad en Flex): emite \texttt{LEXICAL\_ERROR} (999) sin
    fragmentar el lexema, marca \texttt{lexical\_error\_seen} y \textbf{frena}
    el análisis (código de salida~2, sin \texttt{PARSE\_OK}).
  \item \textbf{\texttt{NEWLINE}} devuelve \texttt{TOKEN\_NEWLINE}, terminador
    de sentencia que el parser usa en \texttt{top\_line} y \texttt{block\_stmt}.
\end{itemize}

\lstinputlisting[language=C, firstline=229, lastline=236, caption={Regla \texttt{BAD\_ID}: error léxico atómico con prioridad sobre \texttt{ID}.}]{../lexer/lexico_equipo_11_scanner.l}

El parser consume los tokens sintéticos de indentación en la regla
\texttt{suite}:

\lstinputlisting[language=C, firstline=295, lastline=325, caption={Regla \texttt{suite} en \texttt{parser/parser.y}: consumo de INDENT/DEDENT.}]{../parser/parser.y}

\subsubsection{Sección 3 --- Código C auxiliar (\texttt{\%\%} final)}

Tras el segundo \texttt{\%\%} se definen las funciones que permiten al scanner
asistir al parser:

\begin{itemize}
  \item \texttt{bol\_dedent\_token()}: cuando un token inicia en columna~0,
    con la pila de indentación activa y sin paréntesis/corchetes abiertos,
    emite el \texttt{DEDENT} de cierre de bloque antes del token real.
  \item \texttt{defer\_token(tok)}: guarda el token real (y su
    \texttt{yylval}) para emitirlo \emph{después} del \texttt{INDENT}/%
    \texttt{DEDENT} inyectado; soporta el diferimiento de un token.
  \item \texttt{lexer\_multistmt\_error(tok)} y
    \texttt{lexer\_note\_token(tok)}: registran el historial de tokens
    (omitiendo \texttt{NEWLINE}/\texttt{INDENT}/\texttt{DEDENT}) y detectan
    patrones de \textbf{múltiples sentencias en una línea}
    (\texttt{pass pass}, \texttt{return 5 6}, \texttt{x = 5 y = 6}),
    señalando el error que complementa a \texttt{validate\_stmt\_line\_end}
    del parser.
\end{itemize}

\lstinputlisting[language=C, firstline=362, lastline=409, caption={Funciones auxiliares del scanner que asisten al parser: \texttt{defer\_token}, \texttt{lexer\_multistmt\_error}, \texttt{lexer\_note\_token} y \texttt{bol\_dedent\_token} (\texttt{lexer/lexico\_equipo\_11\_scanner.l}).}]{../lexer/lexico_equipo_11_scanner.l}
